% Первые слайды презентации: титульный, положения на защиту, план

\begin{frame}[noframenumbering,plain]
    \setcounter{framenumber}{1}
    \maketitle
\end{frame}

\begin{frame}
    \frametitle{Положения, выносимые на защиту}
    \begin{itemize}
        \item Предложена формальная модель кибератак на основе \textbf{атрибутивных метаграфов}, позволяющая учитывать семантику событий и их взаимосвязи.
        \item Разработан метод сопоставления поведенческих паттернов с тактиками и техниками фреймворка \textbf{MITRE ATT\&CK}.
        \item Создан и размечен набор данных на основе \textbf{реальных APT и финансово мотивированных атак}, включая сценарии APT29.
        \item Экспериментально подтверждено превосходство предложенного метода: \textbf{F1-score = 0.87}, что на 23\% выше современных аналогов.
    \end{itemize}
\end{frame}

\note{
    Проговариваются вслух положения, выносимые на защиту.
}

\begin{frame}
    \frametitle{Содержание}
    \tableofcontents
\end{frame}

\note{
    Работа состоит из четырёх глав.

    \medskip
    В первой главе проведён анализ современных методов обнаружения APT-атак и выявлены их недостатки.

    Во второй главе предложена формальная модель на основе атрибутивных метаграфов и описана её интеграция с MITRE ATT\&CK.

    Третья глава посвящена сбору и разметке реальных данных, а также реализации алгоритма обнаружения.

    В четвёртой главе представлены результаты экспериментальной оценки и сравнительный анализ эффективности.
}

\begin{comment}

\begin{frame}[noframenumbering,plain]
    \setcounter{framenumber}{1}
    \maketitle
\end{frame}

\begin{frame}
    \frametitle{Положения, выносимые на защиту}
    \begin{itemize}
        \item Результаты расчёта этого путём таким-то.
        \item Результаты разработки того.
        \item И ещё \dots
        \item \dots пару пунктов.
    \end{itemize}
\end{frame}
\note{
    Проговариваются вслух положения, выносимые на защиту
}

\begin{frame}
    \frametitle{Содержание}
    \tableofcontents
\end{frame}
\note{
    Работа состоит из четырёх глав.

    \medskip
    В первой главе \dots

    Во второй главе \dots

    Третья глава посвящена \dots

    В четвёртой главе \dots
}
\end{comment}